\section*{अध्याय \ref{ch:b2:enolates-aldol} --- ईनॉल, ईनोलेट और ऐल्डोल अभिक्रिया}

\begin{solution}{exo:b2:enolates-aldol:1}
ब्यूटेनोन: ब्यूट-1-ईन-2-ऑल \ce{CH2=C(OH)CH2CH3} और ब्यूट-2-ईन-2-ऑल \ce{CH3C(OH)=CHCH3} ((\textit{E}) और
(\textit{Z}))। साइक्लोहेक्सेनोन: साइक्लोहेक्स-1-ईन-1-ऑल।
\end{solution}

\begin{solution}{exo:b2:enolates-aldol:2}
एथेन $<$ प्रोपेनोन $<$ एथिल 3-ऑक्सोब्यूटेनोएट (10.7) $<$ पेंटेन-2,4-डाइओन (9.0)।
\end{solution}

\begin{solution}{exo:b2:enolates-aldol:3}
3-हाइड्रॉक्सीब्यूटेनल, फिर गर्म करने पर ब्यूट-2-ईनैल।
\end{solution}

\begin{solution}{exo:b2:enolates-aldol:4}
डाइएथिल प्रोपेनडाइओएट, सोडियम एथॉक्साइड, 1-ब्रोमोब्यूटेन; ब्यूटिलप्रोपेनडाइओइक अम्ल तक जल-अपघटन;
गर्म करने से \ce{CO2} हटता है: \ce{CH3(CH2)3CH2COOH}, हेक्सेनोइक अम्ल।
\end{solution}

\begin{solution}{exo:b2:enolates-aldol:5}
गतिक: C6 पर (\ce{CH2} वाली ओर) विप्रोटॉनन, \qty{-78}{\degreeCelsius} पर LDA से। ऊष्मागतिक:
अधिक प्रतिस्थापित द्वि-आबंध C1=C2, मेथिल धारण करने वाले कार्बन की ओर, कमरे के ताप पर अपने ऐल्कोहॉल में
किसी ऐल्कॉक्साइड से।
\end{solution}

\begin{solution}{exo:b2:enolates-aldol:6}
4-फ़ेनिलब्यूट-3-ईन-2-ओन (बेन्ज़ैल्ऐसीटोन), \ce{C6H5CH=CHCOCH3}। बेन्ज़ैल्डिहाइड का कोई $\alpha$ हाइड्रोजन नहीं है: वह
\omterm{def:b2:enolates-aldol:enolate}{ईनोलेट} नहीं बना सकता और केवल इलेक्ट्रॉनरागी है।
\end{solution}

\begin{solution}{exo:b2:enolates-aldol:7}
एक मेथिल समूह (C1) का \omterm{def:b2:enolates-aldol:enolate}{ईनोलेट} दूसरे कार्बोनिल (C5) पर आक्रमण करता है: एक पंचसदस्यी वलय,
निर्जलीकरण के बाद 3-मेथिलसाइक्लोपेंट-2-ईन-1-ओन। विकल्प, C3 पर बना \omterm{def:b2:enolates-aldol:enolate}{ईनोलेट} C5 पर आक्रमण करता हुआ,
एक तनावग्रस्त त्रिसदस्यी वलय बनाता।
\end{solution}

\begin{solution}{exo:b2:enolates-aldol:8}
सोडियम एथॉक्साइड के साथ \ce{2CH3COOC2H5 -> CH3COCH2COOC2H5 + C2H5OH}; विप्रोटॉनित उत्पाद
अम्लीकरण पर एथिल 3-ऑक्सोब्यूटेनोएट के रूप में प्राप्त होता है।
\end{solution}

\begin{solution}{exo:b2:enolates-aldol:9}
प्रोपेनोन, एथेनल और ऐसीटोफ़ीनोन (इनमें \ce{CH3-CO} समूह है; एथेनल \ce{CH3CHO} है)।
पेंटेन-3-ओन नहीं देता।
\end{solution}

\begin{solution}{exo:b2:enolates-aldol:10}
एक सिरे का \omterm{def:b2:enolates-aldol:enolate}{ईनोलेट} दूसरे सिरे के एस्टर पर आक्रमण करता है: एथिल 2-ऑक्सोसाइक्लोपेंटेन-1-कार्बोक्सिलेट, एक
पंचसदस्यी वलय।
\end{solution}

\begin{solution}{exo:b2:enolates-aldol:11}
3-मेथिलपेंटेन-2-ओन में त्रिविमजनक केंद्र (C3) ही $\alpha$ कार्बन है: ईनॉलन उसे समतलीय बना देता है, और
किसी भी फलक से पुनःप्रोटॉनन उसे रेसेमिक कर देता है। 4-मेथिलहेक्सेन-2-ओन में त्रिविमजनक केंद्र C4 पर है,
कार्बोनिल के सापेक्ष $\beta$, जिसे ईनॉलन नहीं छूता।
\end{solution}

\begin{solution}{exo:b2:enolates-aldol:12}
क्षारक में बेन्ज़ैल्डिहाइड के दो तुल्यांकों के साथ प्रोपेनोन। दो भिन्न ऐल्डिहाइडों के साथ पहले
प्रोपेनोन को एक ऐल्डिहाइड से संघनित कीजिए (एक तुल्यांक, बेन्ज़ैल्ऐसीटोन को पृथक करते हुए), फिर उसके
शेष मेथिल समूह को दूसरे से संघनित कीजिए। एक ही पात्र में दोनों ऐल्डिहाइड प्रतिस्पर्धा करते हैं और
तीन डाइईनोनों का मिश्रण देते हैं।
\end{solution}

\begin{solution}{pb:b2:enolates-aldol:1}
\textbf{1.} प्रोपेनोन, छह $\alpha$ हाइड्रोजन (हर मेथिल पर तीन)।
\textbf{2.} $\mathrm{CH_3{-}CO{-}CH_2^-} \leftrightarrow \mathrm{CH_3{-}C(O^-){=}CH_2}$।
\textbf{3.} इसके कार्बोनिल के बगल वाले कार्बन पर कोई हाइड्रोजन नहीं है (वह कार्बन वलय का है और
उस पर कोई H नहीं)।
\textbf{4.} बेन्ज़ैल्डिहाइड, आधिक्य में एक ऐल्डिहाइड, बेहतर इलेक्ट्रॉनरागी है; प्रोपेनोन का \omterm{def:b2:enolates-aldol:enolate}{ईनोलेट} उससे
प्रोपेनोन कार्बोनिल की तुलना में अधिक बार मिलता है।
\textbf{5.} कीटोन के दो कार्बन समूह उसके कार्बोनिल कार्बन पर इलेक्ट्रॉन घनत्व धकेलते हैं और
निकट आने में बाधा डालते हैं; ऐल्डिहाइड का एक।
\textbf{6.} \omterm{def:b2:enolates-aldol:enolate}{ईनोलेट} का केवल छोटा अंश चाहिए: अभिक्रिया करते ही उसका पुनर्जनन होता है।
\textbf{7.} \omterm{def:b2:enolates-aldol:enolate}{ईनोलेट} कार्बन बेन्ज़ैल्डिहाइड के कार्बोनिल कार्बन से आबंधित होता है:
4-हाइड्रॉक्सी-4-फ़ेनिलब्यूटेन-2-ओन का ऐल्कॉक्साइड।
\textbf{8.} एक $\alpha$ हाइड्रोजन का निष्कासन और $\beta$ कार्बन से हाइड्रॉक्साइड की हानि: 4-फ़ेनिलब्यूट-3-ईन-2-ओन।
\textbf{9.} बना द्वि-आबंध वलय और कार्बोनिल दोनों से संयुग्मित है।
\textbf{10.} दूसरे मेथिल समूह पर अभी भी तीन $\alpha$ हाइड्रोजन हैं।
\textbf{11.} \ce{2C6H5CHO + CH3COCH3 -> C6H5CH=CHCOCH=CHC6H5 + 2H2O}।
\textbf{12.} निर्जलीकरण, और उत्पाद का अवक्षेपण, जो विलयन छोड़ देता है।
\textbf{13.} दो, दोनों (\textit{E})।
\textbf{14.} (\textit{Z}) समावयवी फ़ेनिल समूह को कार्बोनिल समूह के बगल में रखते हैं: त्रिविम तनाव।
\textbf{15.} दो \ce{C=C}, एक \ce{C=O} और दोनों वलय: पूरे अणु में फैला एक संयुग्मित निकाय।
\textbf{16.} इसका विस्तारित संयुग्मन $\pi$ अंतराल को संकरा करता है (\cref{prop:b2:huckel:gap}): यह दृश्य क्षेत्र के
नीले सिरे पर अवशोषण करता है और पीला दिखता है।
\textbf{17.} नहीं: यह समतलीय (या लगभग) है और इसका एक दर्पण तल है।
\textbf{18.} 106.12, 58.08 और \qty{234.30}{g/mol}।
\textbf{19.} बेन्ज़ैल्डिहाइड $2.1 \times 1.045/106.12 = \qty{20.7}{mmol}$; प्रोपेनोन $0.75 \times
0.7845/58.08 = \qty{10.1}{mmol}$।
\textbf{20.} प्रोपेनोन को बेन्ज़ैल्डिहाइड के \qty{20.3}{mmol} चाहिए, और \qty{20.7}{mmol} उपलब्ध है:
प्रोपेनोन सीमांत है, बस थोड़े अंतर से।
\textbf{21.} $0.01013 \times 234.30 = \qty{2.37}{g}$।
\textbf{22.} बचा हुआ बेन्ज़ैल्डिहाइड एथेनॉल में रहता है और धुल जाता है।
\textbf{23.} एकल-संघनन उत्पाद, बेन्ज़ैल्ऐसीटोन।
\textbf{24.} $0.75 \times 2.37 = \boldsymbol{\approx \qty{1.78}{g}}$ डाइबेन्ज़ैल्ऐसीटोन।
\end{solution}
